\section{An optional Galois-orbit criterion}
\label{sec:galois}
This section gives a sufficient arithmetic condition for a bounded detector locus to fill a full unitary family. It uses Andr\'e--Oort and is separate from the proof of Theorem~\ref{thm:global}.

Fix a geometrically connected fine unitary Shimura base $S$ over a number field $F$, its polarized abelian scheme, and a labelled integral order action in $K$, all over that same $F$. Suppose that $G^{\rm ad}$ is $\Q$-simple of type $A$, as for the full signature-$(n,n)$ unitary datum. A special point is \emph{relative Weyl} if its defining torus has Galois action containing the Weyl group of this datum \cite[Definitions~5.3--5.4]{ChaiOort}. This refers to the unitary group; a full Weyl CM field for the ambient Siegel variety is not required.

Choose a locally closed embedding $S\hookrightarrow\mathbf P^N$. For a reduced closed set $V$ let $\operatorname{tdeg}(\overline V)$ be the sum of the degrees of all geometric irreducible components of its projective closure, in every dimension. This convention matters: a line with arbitrarily many additional isolated points has fixed top-dimensional degree.

Assume Andr\'e--Oort for every Cartesian power of $S$. Scanlon's automatic uniformity theorem \cite[Theorem~4.2]{Scanlon} gives a function $B_S(H)$ bounding the number of special points on a closed locus of total degree at most $H$ that lie on no positive-dimensional special subvariety contained in that locus. To obtain this form, parameterize bounded-degree reduced closed sets by finitely many Chow families. Automatic uniformity bounds the families of Zariski closures of their special points, and hence bounds the numbers of their zero-dimensional components. Andr\'e--Oort identifies all the positive-dimensional components as special. The resulting function is qualitative; no numerical rate is asserted.

\begin{proposition}\label{prop:galois-witness}
Let $\mathcal D_D=\{s:\delta_W(s)\leq D\}$ and choose any integer $H(D)$ at least $\operatorname{tdeg}(\overline{\mathcal D_D})$. Under the preceding hypotheses, a relative-Weyl special point $x$ satisfying
\[
 \delta_W(x)\leq D,\qquad
 |\operatorname{Gal}(\overline F/F)\cdot x|>B_S(H(D))
\]
forces $\mathcal D_D=S$. More generally, if $\mathcal D_D$ is proper, it contains at most $B_S(H(D))$ relative-Weyl special points.
\end{proposition}

\begin{proof}
First, the locus $\mathcal D_D$ descends over the fixed field $F$, independently of $D$. Choose one global order endomorphism $a$ with $a/\bar a$ not a root of unity. The integral operator $P(a^*+\bar a^*)$ in the proof of Theorem~\ref{thm:extraction} is defined over $F$ and detects exactly the Weil projection. On a projective parameter space for effective cycles of degree at most $D$, apply this operator to the universal class in one fixed $\ell$-adic cohomology theory. Its nonzero locus is a finite union of geometric connected components, stable under $\operatorname{Gal}(\overline F/F)$ because the operator and the class construction are defined over $F$. This open and closed subscheme descends over $F$; its proper image is precisely $\mathcal D_D$. Detection by a reducible effective cycle is equivalent to detection by at least one of its components. Thus the entire $F$-orbit of $x$ is contained in $\mathcal D_D$.

A relative-Weyl special point lies on no proper positive-dimensional special subvariety of $S$. Its Mumford--Tate torus maps onto the full adjoint Weyl torus \cite[Remark~5.8]{ChaiOort}. The group of a special subvariety through the point contains that torus. Its roots are stable under Galois and hence under the full Weyl group. In type $A$ the roots of each absolute simple factor form one Weyl orbit; $\Q$-simplicity makes Galois transitive on the factors. The subgroup therefore has either no roots, giving a zero-dimensional domain, or all roots, giving the whole connected domain. This is the calculation in \cite[Lemma~6.3]{ChaiOort}.

The same property holds for Galois conjugates, because Galois conjugates of special subvarieties are special. If $\mathcal D_D$ were proper, every point in the indicated orbit would therefore be one of its isolated maximal special points. The bound $B_S(H(D))$ contradicts the assumed orbit size. The same reasoning counts all relative-Weyl points in a proper $\mathcal D_D$.
\end{proof}

The detector degree $D$, the degree $H(D)$ of its locus in the base, and the special-point bound $B_S(H(D))$ are different quantities. Closedness supplies no growth estimate for $H(D)$. Automatic uniformity supplies no effective estimate for $B_S$. Most importantly, an orbit lower bound supplies no detector membership. A pointwise detector of degree $D(x)$ is useful here only if the comparison with $B_S(H(D(x)))$ is proved.

Without the relative-Weyl restriction, orbit size alone is insufficient. In $Y(1)^2$, the fixed special line $\{0\}\times Y(1)$ contains the full Galois orbits of $(0,j)$ for singular moduli $j$, with unbounded orbit sizes. On a curve base, in contrast, every proper closed locus is zero-dimensional, and its total degree directly bounds the size of an orbit it contains.
